% TOP_PROBLEM: {"id": 3273, "title": "Large acyclic induced subdigraph in a planar oriented graph.", "category_id": 3, "category": {"id": 3, "name": "graph_theory", "display_name": "Graph Theory", "description": "Problems involving graphs, networks, and their properties.", "slug": "graph-theory", "order_index": 3, "created_at": "2026-07-31T15:26:25.671Z"}, "status": "open", "problem_number": "OPG-52197", "set_id": 12}
% FIRST_POSTED: 2026-10-02
% ATTEMPT_STATUS: unresolved
% UnsolvedMath ID 3273; source catalogue code OPG-52197
% !TeX program = lualatex
% GPT-6 Astra Ultra research attempt; independently checked partial work, 2026-10-02.
\documentclass[11pt,a4paper]{article}
\usepackage[margin=25mm]{geometry}
\usepackage{fontspec}
\setmainfont{lmroman10-regular.otf}[BoldFont=lmroman10-bold.otf,
 ItalicFont=lmroman10-italic.otf,BoldItalicFont=lmroman10-bolditalic.otf]
\usepackage{amsmath,amssymb,mathtools}
\usepackage{unicode-math}
\setmathfont{latinmodern-math.otf}
\AtBeginDocument{\let\setminus\smallsetminus}
\usepackage{xurl,hyperref}
\hypersetup{colorlinks=true,urlcolor=blue}
\setcounter{secnumdepth}{0}
\setlength{\parindent}{0pt}
\setlength{\parskip}{5pt}
\setlength{\emergencystretch}{3em}
\begin{document}
\section*{Large acyclic induced subdigraph in a planar oriented graph.}\label{3273}
{\small UnsolvedMath ID 3273 · OPG-52197 · Graph Theory · OpenGarden}
\subsection{Definitions and mathematical statement}
Let $D=(V,A)$ be a finite oriented graph whose underlying simple
undirected graph is planar. Thus there are no loops, repeated arcs,
or opposite pairs of arcs. Put $n=|V|$. For $U\subseteq V$, the induced
subdigraph $D[U]$ contains all arcs of $D$ having both ends in $U$.
It is acyclic when it has no directed cycle, equivalently when its
vertices admit a linear ordering in which every arc points forward.

The conjecture asks whether one can always find $U\subseteq V$ with
\[
                           |U|\ge\frac35n
                  \quad\text{and}\quad D[U]\text{ acyclic}.
\]
Since $|U|$ is an integer, the size condition means
$|U|\ge\lceil3n/5\rceil$. The equivalent deletion formulation asks for
a set $S\subseteq V$ of at most $\lfloor2n/5\rfloor$ vertices meeting
every directed cycle; then $U=V\setminus S$. Such an $S$ is a directed
feedback vertex set.

Deleting arcs is not allowed when testing whether $D[U]$ is acyclic:
all arcs between retained vertices remain. On the other hand,
undirected cycles are allowed if their orientations do not form
directed cycles. This distinguishes the target from an induced-forest
problem for the underlying graph. It also asks for one large acyclic
set, not a partition of all vertices into a fixed number of acyclic
classes. The graph may be disconnected, and the orientation is arbitrary
subject to planarity of the underlying graph.
\subsection{Short English statement}
Can at least three fifths of the vertices of every planar oriented graph be retained without retaining any directed cycle?
\subsection{Research attempt: verification of a known negative answer}
\paragraph{Scope and literature correction.}
GPT-6 Astra Ultra, 2 October 2026. The historical assertion transcribed
above is false. Knauer, Valicov and Wenger [S3] constructed planar
oriented graphs contradicting it; their result predates this attempt.
The catalogue metadata is preserved for identity and is not evidence
of current mathematical openness. This document verifies a concrete
seven-vertex negative certificate independently, including planarity
and the exact maximum induced acyclic-set size. No new solution
priority is claimed, and no unproved extension is needed to refute
the displayed universal assertion.

\paragraph{1. Explicit oriented graph.}
Use vertices $0,1,\ldots,6$ and precisely the following thirteen arcs:
\[
\begin{split}
 &0\to1,\quad1\to2,\quad2\to0,\\
 &3\to4,\quad0\to3,\quad1\to3,\quad4\to0,\quad4\to1,\\
 &5\to6,\quad3\to5,\quad4\to5,\quad6\to3,\quad6\to4.
\end{split}
\]
There are no loops and no opposite pairs. Since these are the full
arc data, the induced-subgraph requirement is unambiguous: no arc
may subsequently be omitted when deciding whether retained vertices
are acyclic.

\paragraph{2. A constructive planarity certificate.}
Begin with the undirected triangle on $0,1,2$ drawn in the plane.
Insert vertex $4$ inside its bounded triangular face and join it to
all three boundary vertices. Next insert vertex $3$ inside face
$014$, joining it to $0,1,4$. Insert vertex $6$ inside face $034$,
joining it to $0,3,4$. Finally insert vertex $5$ inside face $346$,
joining it to $3,4,6$. Each insertion takes place in an existing
triangular face and divides it into three triangular faces, so it
introduces no crossing.

The resulting planar supergraph has fifteen edges. Deleting edges
$24$ and $06$ leaves exactly the underlying graph of the thirteen
arcs above. Planarity is preserved under these deletions. This
certificate does not rely on a drawing inferred from an arc list
or on a black-box planarity routine; the indicated face sequence
constructs an embedding directly.

\paragraph{3. Every feedback vertex set has at least three vertices.}
The graph contains directed triangles with vertex sets
\[
 \{0,1,2\},\quad\{0,3,4\},\quad\{1,3,4\},\quad
 \{3,5,6\},\quad\{4,5,6\}.
\]
Their cyclic directions are read directly from the arc list.
Suppose a feedback vertex set $S$ had at most two vertices.
If it contains $5$ or $6$, at most one other vertex remains
available to hit the first three triangles. Those three vertex
sets have empty common intersection, so one vertex cannot hit
them all. This is impossible.

If $S$ contains neither $5$ nor $6$, hitting the last two triangles
forces both $3$ and $4$ to lie in $S$. But then the first triangle
$012$ remains. This is also impossible. Thus every feedback vertex
set has at least three vertices, and every induced acyclic set has
size at most $7-3=4$. The argument uses only necessary hits on five
specified cycles; additional cycles cannot invalidate the upper bound.

\paragraph{4. Matching lower bound and the failed threshold.}
Retain vertices $\{1,2,4,6\}$. The complete induced arc set is
$6\to4$, $4\to1$, and $1\to2$, a directed path. These four vertices
therefore form an induced acyclic set, proving that the maximum size
is exactly four. The historical claim would require at least
$\lceil3\cdot7/5\rceil=5$ vertices, which is impossible.

Taking disjoint unions of $q$ copies gives planar oriented graphs
on $7q$ vertices with maximum acyclic-set size exactly $4q$.
The maximum is additive over components because a directed cycle
cannot cross components. Since $4q<21q/5$, the contradiction is
not restricted to one small order or a rounding artefact.

\paragraph{Verification and mathematical status.}
The script [S9] verifies each face insertion, the exact two deleted
edges, all five triangle certificates, and all $128$ vertex subsets.
It confirms maximum acyclic-set size four. The negative answer is
already established in [S3]; the present contribution is a fully
specified and reproducible certificate. There is no remaining gap
in refuting the historical $3n/5$ assertion. This does not settle
any stronger lower-bound question such as the separate half-order
problem discussed in the literature.

\paragraph{Final status.}
Historical target: DISPROVED in the cited literature, with its negative
answer verified here. The conservative repository attempt flag is
UNRESOLVED as a workflow classification; it is not a claim that this
historical assertion remains mathematically open or a claim of a new
resolution by this model.
\subsection{Sources}
\begin{itemize}
\item[S1] ULAM.AI and the UnsolvedMath Contributors, supplied problems.json, record 3273 (OPG-52197), OpenGarden collection. Catalogue identity and source transcription; CC BY 4.0.
\url{https://huggingface.co/datasets/ulamai/UnsolvedMath}.
\item[S2] Open Problem Garden, Large acyclic induced subdigraph in a planar oriented graph.. Original problem and discussion; the mathematical formulation below is expanded from the supplied source record.
\url{http://www.openproblemgarden.org/op/large_acyclic_induced_subdigraph_in_a_planar_oriented_graph}.
\item[S3] Kolja Knauer, Petru Valicov and Paul S. Wenger, \emph{Planar digraphs without large acyclic sets} (2015; revised 2016), Journal of Graph Theory. \url{https://arxiv.org/abs/1504.06726}; \url{https://doi.org/10.1002/jgt.22061}.
\item[S9] GPT-6 Astra Ultra, exact finite checks accompanying this attempt (2 October 2026). Repository script, Python standard library only. \texttt{scripts/check\_attempt\_batch03\_digraphs\_2026\_10\_02.py}.
\end{itemize}
\end{document}
